\documentclass[11pt,a4paper]{article}

\usepackage[utf8]{inputenc}
\usepackage[cyr]{aeguill}
\usepackage[francais]{babel}

% Les packages
%\usepackage[french]{babel}
\usepackage{fancybox}
\usepackage{graphics}
\usepackage{color}
\usepackage{latexsym}
\usepackage{amsmath,amsthm,amssymb}
\usepackage[plain]{fullpage} 
\usepackage{verbatim}
\usepackage{times,epsfig}
\usepackage{listings}
\usepackage{multirow}
\usepackage{amsmath,amsthm,amssymb,amscd,txfonts,bbm} % RAJOUT ICI !!!
\usepackage{graphics,epsfig} % RAJOUT ICI !!!
%\usepackage{algorithmic}
\usepackage[final]{pdfpages} 
\usepackage{fancyvrb}
\usepackage{enumerate}
\usepackage{listingsutf8}
\usepackage{comment}
\usepackage{minted}

% small et verbatim: to be coherent with \file
\newenvironment{smallverbatim}%
{\endgraf\small\verbatim}{\endverbatim}

\newenvironment{qv}
{\quote\Verbatim}
{\endquote\endVerbatim}

\usepackage{lastpage}

\newcommand{\terme}[1]{\texttt{#1}}
\newcommand{\termSet}[1]{\{\texttt{#1}\}}
\newcommand{\guill}[1]{<<\,#1\,>>}
% Macro pour l'inclusion de figure  
\newcommand{\fig}[2]{%
  \begin{figure*}[hbt]
   \begin{center}
      \includegraphics[width=12cm]{#1}
  \caption{\label{#1}  #2}
  \end{center}
 \end{figure*}
}
\newcommand{\figTex}[2]{%
  \begin{figure*}[hbt]
 \centerline{\input{#1.pstex_t}}
  \caption{\label{#1}  #2}
 \end{figure*}
}

% Macro résumés
\definecolor{turquoise}{rgb}{0.20 ,0.60, 0.60}
\definecolor{vert}{rgb}{0.40,  0.60 ,0.00}
\definecolor{violet}{rgb}{0.80 , 0.00 , 0.80}
\definecolor{bleu}{rgb}{0.20,  0.20,  0.80}
\definecolor{rose}{rgb}{1.00,  0.20,  0.40}

\lstset{basicstyle=\footnotesize,showstringspaces=false,language=XML,breaklines=false,columns=fullflexible,frame=shadowbox, captionpos=b}

\newenvironment{solution}[0]{
~\\ \color{red} \textbf{Solution : }\color{black}}{%
}
%\excludecomment{solution}


\newcommand{\cle}[1]{\textbf{#1}}
\def\fk#1{\textit{#1}}
\def\tble#1{\textit{#1}}

\parindent=0pt           
                             
\usepackage{url}

\newtheorem{Exercice}{Definition}

\begin{document}

\title{Sujet d'examen UE NFE204 - Session 2}
\author{Année universitaire 2020-2021}
\date{Examen seconde session: 31 août 2021}

\maketitle

{\bf \Large Consignes.}
Cet examen est exceptionnellement organisé à distance. 
Vous devez rendre votre copie numérisée à 21h15 au plus tard. 
Aucune copie ne sera acceptée après cette heure limite. 
Les seuls formats acceptés sont le PDF et le JPEG, ceci afin d'éviter tout problème de lecture.  
Prenez vos dispositions pour être en mesure de fournir le bon format le moment venu. Une activité 
de discussion en direct (\guill{\textit{chat}}) est ouverte dans l'espace Moodle pour d'éventuelles questions.


\section*{Première partie : modélisation semi-structurée (8 pts)}

Un grand magasin enregistre tous les tickets de caisse des transactions effectuées par ses clients.
Un ticket de caisse comprend la date et l'heure de chaque transaction, le numéro de la caisse, le prix total,
et la liste des articles achetés avec, pour chacun, la quantité et le prix unitaire.  Le ticket de caisse comprend également
l'identifiant du client si ce dernier a utilisé sa carte de fidélité. Sinon cet identifiant
est à \texttt{null}.


\begin{enumerate}
  \item Dans une base relationnelle, combien de tables faudrait-il pour représenter 
         les informations d'un ticket de caisse (expliquer et donner le schéma de ces tables); combien d'insertions de ligne
             faut-il faire pour chaque ticket\,?
    \item Si on choisit une base NoSQL, comment représenter un ticket de caisse en JSON.
        Donnez un exemple simple. Combien d'insertions faut-il faire\,?
   \item Je veux compter le nombre de transactions par caisse et par jour.  Quelle est la requête SQL si les données
       sont dans une base relationnelle\,?
    \item Même question si la base est NoSQL\,: comment calculer le nombre de transactions par caisse et par jour
    \item Proposez un format de document semi-structuré (en JSON par exemple) pour une base
         des articles, la description de chaque article incluant la liste des ventes (une vente = date,
         quantité, client)
 \end{enumerate}


\begin{solution}

\begin{enumerate}
   \item En relationnel, il faut 2 tables, une pour les tickets avec les champs \texttt{noCaisse},
      \texttt{date}, \texttt{heure} et \texttt{prixTotal}, une autre pour les articles avec \texttt{idArticle},
         \texttt{quantité} et \texttt{prixUnitaire}. Un champ \texttt{idClient}
          dans la table \texttt{Ticket} peut être à \texttt{null}. La table \texttt{Article} comprend
         une clé étrangère \texttt{idTicket} pour savoir de quel ticket l'achat de l'article fait partie.
         
          Il faut une insertion dans la table \texttt{Ticket},
           et autant d'insertions dans la table \texttt{Article} qu'il y a d'articles achetés dans
           une même transaction. 
      
      \item En JSON, par exemple:

\begin{minted}[frame=leftline,
               baselinestretch=.9,
               fontsize=\footnotesize,
               framesep=2mm]{json}
{
  "_id": 978,
  "noCaisse": 12,
  "date": "01/01/2017",
  "heure": "13:30",
  "prixTotal": 58.47,
  "idClient": null,
  "articles": [
     {"idArticle": 1982, "quantité": 5,   "prix": 10.45},
     {"idArticle": 34, "quantité": 1,   "prix": 45},
     {"idArticle": 154, "quantité": 3,   "prix": 3.02},
  ]
}
\end{minted}

Une insertion suffit.

\item En SQL:

\begin{minted}{sql}
select count(*) from Ticket t, Article a 
where a.idTicket = t.id
group by date, noCaisse
\end{minted}

\item En NoSQL, il faut écrire et exécuter un programme MapReduce.

La fonction de Map émet  une paire intermédiaire dont la clé est le critère de regroupement,
soit la paire (date, numéro de caisse).

\begin{minted}{bash}
function fonctionMap ($ticket)
{
  emit ( [$ticket.date, $ticket.noCaisse], 1)
}
\end{minted}

La fonction de Reduce effectue le compte des paires intermédiaires, autrement dit des tickets
ayant même  date et même numéro de caisse.

\begin{minted}{bash}
function fonctionReduce ($clé, $valeurs)
{
  return ($clé, count($valeurs))
}
\end{minted}
\end{enumerate}
\end{solution}


\section*{Deuxième partie : MapReduce (7 pts)}

Je veux calculer le montant moyen des tickets par client, pour les clients dont
je connais l'identifiant. 

\begin{enumerate}
 \item Décrivez la modélisation MapReduce de ce calcul. 
 Donnez le pseudo-code de chaque
fonction, ou indiquez par un texte clair son déroulement.
\item Donnez la forme de la chaîne de traitement (\textit{workflow}), avec des opérateurs Pig ou Spark,
      qui implante ce calcul.
   \item Indiquez comment on produit un document 'Article' avec la liste des ventes à partir de la
    collection des tickets de caisse.
\end{enumerate}


\begin{solution}

\begin{enumerate}
 
\item La fonction de Map  ne prend en compte que les tickets ayant un numéro de client non nul,
   et émet  une paire intermédiaire dont la clé est le critère de regroupement,
soit le numéro de client, et la valeur le prix total du ticket.

\begin{minted}{bash}
function fonctionMap ($ticket)
{
  if ($ticket.idClient not null) then
      emit ( $ticket.idClient, $ticket.prixTotal)
  end if
}
\end{minted}

La fonction de Reduce effectue la moyenne des prix.
\begin{minted}{bash}
function fonctionReduce ($idClient, $listePrix)
{
  return ($idClient, avg($listePrix))
}
\end{minted}

\item En Pig:

\begin{minted}{pig}
tickets_client = filter tickets BY idClient not null;
tickets_groupes = group tickets_client by idClient;
avg_par_client = foreach tickets_groupes generate group, AVG(tickets_groupes.prixTotal);
\end{minted}
\end{enumerate}
\end{solution}

\section*{Troisème partie : systèmes distribués (3 pts)}

On vous demande d'expliquer quelques-uns des principes du \textit{consistent hashing}. 

\begin{enumerate}
   \item Expliquez comment on choisit le serveur sur lequel on place un document.
  \item Pourquoi la structure est-elle basée sur un anneau? Pourquoi pas un segment de droite par exemple?
  \item Pourquoi y a-t-il autant de valeurs sur l'anneau (2 puissance 64 typiquement)? Pourquoi ne pas commencer
      avec 8 ou 16 valeurs, et augmenter ensuite?
   \item Quelle est l'utilité de la notion de n{\oe}ud virtuel?
\end{enumerate}

\section*{Quatrième partie : questions de cours (2 pts)}

 Pourquoi faut-il stocker les paires intermédiaires sur disque dans un processus
 MapReduce? Pourquoi ne peut-on directement les transmettre aux machines qui effectuent
 la phase de Reduce?
 
\end{document}
